\section{Casos de Evaluation y reproducción}

\subsection{Caso típico de Evaluation}
La evaluation list que enumera el docente incluye tareas como replicating Apple iPad Air site, descripción de video y operación de GUI; todas las tareas se dividen en tres subagent: visual, de texto y de herramientas, y luego regresan al Allcastrater, donde se califican con GRM. En el timeline del guion (por ejemplo, 00:20:50) se observa que, después de que el agent visual extrae la información de la interfaz, el agent de texto agrega de inmediato la explicación y el agent de herramientas ejecuta la operación con comandos, formando un flujo completo de Visual Agency.

\begin{knowledgebox}{Enfoque de las tareas de Evaluation}
\begin{itemize}
    \item Reproducir la landing page de Apple iPad Air: el agent visual captura la interfaz, el agent de herramientas genera CSS/HTML y el agent de texto ofrece la explicación.
    \item Benchmark de descripción de video: varios agent extraen frame, generan narrative y alinean timestamps en paralelo dentro de un single loop.
    \item Operación de GUI: VRL y los GUI agents rastrean en conjunto critical steps como clics y deslizamientos, para asegurar que el tool trace sea reproducible.
\end{itemize}
\end{knowledgebox}

\subsection{Fotogramas clave y registro de slides}
Esta clase no tiene slides, así que toda la evidencia visual proviene de los fotogramas del video. Usamos como figure `cover.jpg`, las capturas tomadas en el segmento clave (00:32:00~00:32:48) y el marco de la demo de agent swarm, y en el caption se anota con claridad el rango de tiempo. Para facilitar el audit posterior, también se preparó una evidence table que registra cada benchmark junto con su timecode correspondiente y la evidencia visual/textual.

\begin{table}[H]
\footnotesize
\centering
\begin{tabular}{p{3cm}p{4cm}p{5cm}}
\toprule
Tarea & Periodo & Evidencia \\
\midrule
Replicate Apple iPad Air & 00:20:50--00:22:10 & El docente muestra cómo descomponer el layout y coloca en la misma página de notes la captura del agent visual junto con el resultado de generación de HTML. \\
Agent swarm efficiency & 00:35:20--00:36:50 & Se ralentizan varios fotogramas para mostrar el resumen de puntajes de \texttt{S mean + max} y la GRM reward line, lo que demuestra una efficiency 4.5 veces mayor. \\
Critical steps gating & 00:41:10--00:42:30 & El control panel del GUI agent y el tool trace registran el consumo de critical steps, lo que explica el control del bandwidth. \\
\bottomrule
\end{tabular}
\caption{Key evaluation tasks y su timecode/evidence correspondiente, para que el audit later confirme que cada segmento estructural usa fotogramas de video.}
\end{table}

\subsection{Lista de verificación para la reproducción}
Para reproducir cada evaluation case, organizamos el material disponible en un checklist con cuatro líneas: Data + Algorithm + Agent + Infra; cada línea debe aportar quote, figure, box y summary table, y al final se usa `summary table` para reunir los hallazgos. Cada figure se marca con \verb|\protect\footnotemark| y se anota el timecode específico, para asegurar que las notas redactadas puedan rastrearse en el audit.

\begin{importantbox}{Checklist de reproducción de las notas}
\begin{enumerate}
    \item Reescribir los metadata (autor, fecha, duración, plataforma) y confirmar que no haya placeholder.
    \item Extraer párrafos de \texttt{subs\_clean}, ordenarlos según la lógica pedagógica e insertar boxes y quotes.
    \item En cada section, indicar el timecode de la figure o el número de página de las slides, y escribir una summary table para asegurar que la evidence sea verificable.
    \item Ejecutar xelatex en dos rondas y confirmar que el PDF page count y la cantidad de boxes cumplan con lo requerido.
\end{enumerate}
\end{importantbox}

\subsection{Resumen de la sección}
Los casos de Evaluation subrayan el flujo `critical tasks -> timecode -> evidence`; basta con reproducir cada benchmark siguiendo el checklist e incorporar la evidence table, los boxes y las figures en las notas para que el audit se apruebe con rapidez.
